\documentclass[12pt]{uz-exame} 
\usepackage{amssymb,amsfonts,amsmath}
\newcommand{\rem}[1]{}
\begin{document} 
\examtitle{BSc Honours in Mathematics Level 1}
\examdate{August 2017}
\hours{2}
\lastA{4}
\rubric{\science}
\lastpage{3}
\coursename{CALCULUS 1}
\coursenumber{HMTH101}
%\setcounter{exsec}{17}
\makeheading
\sectionAscience
\begin{question}
\begin{itemize}
\item[(a)] Solve the following inequality $4|x|<7x-6$.
\item[(b)] Show that $|a+b|\leq |a|+|b|$ for all $a,b\in \mathbb{R}$.
\end{itemize}\marks{4,4}
\end{question}
\begin{question}
\begin{itemize}
\item[(a)] Define what is meant by saying that $l$ is the \textbf{limit} of a sequence $\{u_n\}$. \marks{2}
\item[(b)] Evaluate each of the following limits.\\
(i)\, $\displaystyle \lim _{n\rightarrow \infty}\left(\dfrac{\sqrt{3n^2-5n+4}}{2n-20}\right)$,\quad  \quad (ii)\, $\displaystyle \lim _{n\rightarrow \infty} (n-\sqrt{n^2+1})$. \marks{3,3} 
\end{itemize}
\end{question}
\begin{question}
\begin{itemize}
\item[(a)] Estimate the value of $\sqrt[5]{36}$ using differentials. \marks{4}
\item[(b)] Use the Mean Value Theorem to prove that $\ln(1+x)<x$ for $x>0$.\marks{5}
\end{itemize}
\end{question}
\begin{question}
\begin{itemize}
\item[(a)] Show that $\displaystyle\int x^m\ln x\, dx =x^{m+1}\left(\dfrac{\ln x}{m+1}-\dfrac{1}{(m+1)^2}\right), m\neq -1$. \marks{6}
\item[(b)] Find each of the following integrals.\\
(i)\, $\displaystyle \int \cos ^5 x\, dx$ \quad \quad \quad (ii)\, $\displaystyle \int \dfrac{x}{x+2}\, dx$. \marks{5,4}
\end{itemize}
\end{question}
\newpage
\sectionBscience
\begin{question}
\begin{itemize}
\item[(a)] \begin{itemize}
\item[(i)] Prove that if $0<a<b$, then $a^2<b^2$. \marks{4}
\item[(ii)] If $a^2<b^2$, does it necessarily imply that $a<b$? If so, prove it and if not give an example that illustrates your answer. \marks{4}
\end{itemize}
\item[(b)] \begin{itemize}
\item[(i)] Define in terms of $\varepsilon$ and $\delta$, what is meant by saying that a function $f(x)$ is continuous at a point $x=x_0$. \marks{2}
\item[(ii)] Let $f:\mathbb{R}\to \mathbb{R}$  be the function defined by 
\[f(x) = \left\{ 
\begin{array}{l l}
  x^2+1, & \quad x \geq 2\\
  ax+b, & \quad -1 <x < 2\\
  5-2x, & \quad x \leq -1,\\ \end{array} \right. \]
  where $a$ and $b$ are constants. If $f$ is continuous on $\mathbb{R}$, find the values of $a$ and $b$. \marks{6}
\end{itemize}
\item[(c)]\begin{itemize}
\item[(i)] State the Mean Value theorem. \marks{4}
\item[(ii)] Use the Mean Value Theorem or otherwise to prove that 
\[1-\dfrac{a}{b}<\ln \left(\dfrac{b}{a}\right)<\dfrac{b}{a}-1\] if $0<a<b$. \marks{6}
\item[(iii)] Hence or otherwise show that 
\[\dfrac{1}{6}<\ln (1.2)<\dfrac{1}{5}.\] \marks{4}
\end{itemize}
\end{itemize}
\end{question}
\begin{question}
\begin{itemize}
\item[(a)]\begin{itemize}
\item[(i)] Define what is meant by $\displaystyle \lim _{n\rightarrow \infty} u_n=\infty$. 
\item[(ii)] Show using the definition, that $\displaystyle \lim _{n\rightarrow \infty}(3n+7) =\infty$. 
\item[(iii)] Prove that $\displaystyle \lim _{n\rightarrow \infty}\dfrac{c}{n^p}=0$, where $c\neq 0, p>0$. \marks{3,4,4}
\end{itemize}
\item[(b)] \begin{itemize}
\item[(i)] What is the relationship between differentiability and continuity? 
\item[(ii)] Hence prove this relationship.\marks{3,6}
\end{itemize}
\item[(c)]\begin{itemize}
\item[(i)] If $e^{xy}+y\ln x =\sin 2x$, find $\dfrac{dy}{dx}$. 
\item[(ii)] If $y=x^x, \, x>0$, find $\dfrac{dy}{dx}$. \marks{6,4}
\end{itemize}
\end{itemize}
\end{question}
\newpage
\begin{question}
\begin{itemize}
\item[(a)] Evaluate each of the following limits.\\
(i)\, $\displaystyle \lim _{x\rightarrow 0}\dfrac{x-\sin x}{x^3}$ \quad \quad (ii)\, $\displaystyle \lim _{x\rightarrow 0}x^{\sin x}$. \marks{4,6}
\item[(b)] \begin{itemize}
\item[(i)] Prove that if $f(x)$ is continuous on $[a,b]$, differentiable on $(a,b)$ and\\ if $f'(x)\geq 0$ for every $x$ in $(a,b)$, then $f(x)$ is monotonic increasing on $[a,b]$.
\item[(ii)] Verify Rolle's Theorem for $f(x)=x^2(1-x)^2, \, 0\leq x \leq 1$. \marks{4,6}
\end{itemize}
\item[(c)] Show that the function $f(x)=\dfrac{ax+b}{cx+d}$ is one-to-one or injective if $ad-bc\neq 0$. Find its inverse stating the domains of both the function and its inverse. What happens if $ad-bc=0$? \marks{10}
\end{itemize}
\end{question}
\begin{question}
\begin{itemize}
\item[(a)] Evaluate each of the following integrals.\\
(i)\, $\displaystyle \int \dfrac{\sqrt{x}}{\sqrt{x}+1}\, dx$ \quad \quad (ii)\, $\displaystyle \int \sin ^{-1} x\, dx$. \marks{5,5}
\item[(b)] Find the area of the region bounded by the graph of $f(x)=1-x^2, \, x=0, \, x=1$ and the $x-$axis by calculating the limit of the Riemann Sums. \marks{10}
\item[(c)] Find 
\[\int \dfrac{dx}{(x^2+1)(x+1)}.\] \marks{10}
\end{itemize}
\end{question}
\endexam
\end{document} 